\documentclass[12pt]{report}
\input{../bayesuvius.sty}
\begin{document}

%\chapter{Renormalization Group for PDEs}
%\label{ch-renorm-group-pde}
%
secular define
not bound by time but extends forever
not bound by religious canon

$\vec{r}=(x,y,z)$

$X=(t, \vec{r})$

$\nabla = (\partial_{x}, \partial_{y}, \partial_{z})$ and $\nabla^2 = \nabla\cdot\nabla=\partial^2_x + \partial^2_y + \partial^2_z$.

$\eta=$ concentration,
number of particles per $d$-dim volume



Fick's law

\beq
\vec{J}= -D\nabla \eta
\eeq
Continuity equation

\beq
\partial_t \eta = -\nabla\cdot \vec{J}
\eeq
Diffusion Equation

\beq
\partial_t \eta = \nabla 
\left(
D\nabla \eta
\right)
\eeq

If $D$ is constant
\beq
\partial_t\eta=D\nabla^2 \eta
\eeq
Henceforth, 1 space dimension.
$\nabla^2 =\partial_x^2$

$
\kappa  = 2D
$

\section{DE without long term memory}

\beq
P(x; \s^2)=\frac{1}{\sqrt{2\pi \s}}
\exp\left(\frac{-x^2}{2\s^2}\right)
\eeq

Solve 

\beq
\partial_t \eta = \frac{\kappa}{2}\partial_x^2 \eta
\eeq
for free space (no boundaries, $x\in \RR$)
with initial condition

\beq
\eta(x, 0)=\eta_\rvx(x)=
NP(x; \ell^2)
\eeq

\beq
\eta_\rvx(x)\rarrow N\delta(x)\quad \text{as $\ell\rarrow 0$}
\eeq

\beq 
N(t)=\int_{-\infty}^{\infty}dx\; \eta(x, t)
\eeq

\beq
\partial_t N(t) = \frac{\kappa}{2}\int_{-\infty}^\infty \partial_x^2\eta(x, t)=\frac{\kappa}{2} [\partial_x\eta(x,t)]_{x=-\infty}^{x=\infty}=0
\eeq
$N$ is area under $\eta$ curve, it doesn't change with time.

\beq
\eta(x,t)=
NP(x;\kappa t + \ell^2)
\eeq
for $t\gg \ell^2/\kappa$ (i.e., 
either after a long time or for short Gaussian widths $\ell$),

\beq
\eta(x,t)\sim
NP(x; \kappa t)\quad
\text{as $t\rarrow \infty$ or $\ell\rarrow 0$}
\eeq
We call this an {\bf intermediate times}
approximation, because for truly infinite times, $\eta(x,t)\rarrow \infty$

\beq
\engdim{x}=L, \quad 
\engdim{t}=T, \quad 
\engdim{\eta}= 1/L,\quad 
\engdim{D}= L^2/T,\quad 
\eeq


\beq
\Pi =\eta \sqrt{Dt} ,\quad
\Pi_1 = \frac{x}{\sqrt{D t}}
\eeq

\beq
\Pi = f(\Pi_1)\implies \eta = \frac{1}{\sqrt{Dt}}f\left(\frac{x}{\sqrt{Dt}}\right)
\eeq



\beq
x \to x/b,\quad t \to b^{y_t} t,\quad \eta \to b^{y_\eta} \eta
\eeq

\beq
b^{y_\eta -y_t}\partial_t \eta =
b^{y_\eta+2}\partial^2_x\eta
\eeq

\beq
y_t=-2
\eeq

\beq
\eta(x, t)= b^{-y_\eta}\eta\left(\frac{x}{b}, \frac{t}{b^2}\right)
\eeq
choose  $b=\sqrt{t}$
\beq
\eta(x, t)=b^{-y_\eta}\eta\left(\frac{x}{\sqrt{t}}\right)
\eeq




\section{DE with discontinuous diffusion coefficient}


\beq
\partial_t \eta=
\partial_x\left(D(x)\partial_x \eta\right)
\eeq

\beq
D(x)
= 
\left\{
\begin{array}{ll}
\frac{\kappa}{2}(1+g) & \text{if $x>0$} 
\\
\frac{\kappa}{2} & \text{if $x<0$}
\end{array}
\right\} = \frac{\kappa}{2}(1 + g\indi_{x>0})
\eeq

transmission (interface) continuity conditions

concentration continuity
\beq
\eta(0-, t)=\eta(0+, t)
\eeq

flux continuity
\beq
D_{-}\partial_x \eta(0-, t) = D_{+}\eta(0+, t)
\eeq

Note, this is not
a boundary value problem in the usual sense.
There is no boundary at $x=0$. if there were
\begin{itemize}
\item Dirichlet

\beq \eta(0, t)= \eta_\rvt(x)
\eeq

\item
von Neumann
\beq
\partial_x\eta(x,t)= f(x,t)
\eeq
For example, if there were an impermeable wall at $x=0$,

\beq
\partial_x\eta(x,t)=0
\eeq


\item Robin

\beq
[a \partial_x\eta + b \eta](x, t)=g(x,t)
\eeq
\end{itemize}


\beq
\eta=\eta_0+g \eta_1+ g^2 \eta_2\cdots
\eeq

Zeroth order:

\beq
(\partial_t-\frac{\kappa}{2}\partial_x^2)\eta_0=0
\eeq



\beq
\eta_0(x,t)
=NP(x; \kappa t)
\eeq


First order equation:

\beq
(\partial_t-\frac{\kappa}{2}\partial_x^2)\eta_1
=g
\frac{\kappa}{2}
\partial_x\left(
\indi_{x>0}\partial_x \eta_0
\right)
\eeq

\beq
\eta_1(x,t)
=\frac{\kappa}{2}
\int_0^t d\tau
\int_{-\infty}^\infty d\chi\;
G(x-\chi,t-\tau)
\partial_{\chi}\left[
\indi_{\chi>0}
\partial_{\chi}\eta_0(\chi,\tau)
\right]
\eeq
Integrate by parts, assume $\chi=\pm \infty$  terms vanish

\beq
\eta_1(x,t)
=\frac{\kappa}{2}
\int_0^t d\tau
\int_{0}^\infty d\chi\;
[-\partial_\chi  G(x-\chi,t-\tau)
]
\partial_{\chi}\eta_0(\chi,\tau)
\eeq

\beq
G(x-\chi, t-\tau)= P(x-\chi; \kappa(t-\tau))
\eeq


%\beq
%\partial_\chi\left(
%\indi_{\chi>0}\partial_\chi \eta_0
%\right)
%=
%\indi_{\chi>0}\partial_\chi^2 \eta_0
%+
%\delta(\chi)\partial_\chi \eta_0(0^+,\tau)
%\eeq


\beq
-\partial_\chi G(x-\chi, t-\tau)=\frac{(x-\chi)}{\kappa(t-\tau)}
P(x-\chi; \kappa(t-\tau))
\eeq

\beq
\partial_\chi \eta_0(\chi,\tau)
=
-\;\frac{\chi}{\kappa \tau}N P(\chi; \kappa \tau)
\eeq

Assume $t\gg \tau$  and  $x\gg \chi$ in $P()$ only so 

\beqa
\eta_1(x,t)&=& -\;\frac{\kappa}{2}N
P(x; \kappa t) \int_0^t d\tau
\int_0^\infty d\chi\;\frac{\chi (x-\chi)}{\kappa^2 \tau (t-\tau)}P(\chi;\kappa \tau)
\eeqa

\beq
\hat{\chi}= \chi/\sqrt{\kappa \tau},\quad \hat{x}= x/\sqrt{\kappa \tau}
\eeq

\beq
I(\tau)= \frac{1}{\sqrt{2\pi\kappa \tau}}\int_0^\infty d\hat{\chi}\; \hat{\chi}(\hat{x}-\hat{\chi})\exp(-\hat{\chi}^2/2)
\eeq
\beq
\int_0^\infty dy \; y\exp(-y^2/2)=1,\quad
\int_0^\infty dy \; y^2 \exp(-y^2/2)=\sqrt{2}\;\Gamma(3/2)=\sqrt{\frac{\pi}{2}}
\eeq

\beq
I(\tau)= \frac{1}{\sqrt{2\pi \kappa \tau}}\left[\hat{x}-\sqrt{\frac{\pi}{2}}\right]
\eeq


\beqa
\eta_1(x,t)&=& -\frac{1}{2}
NP(x; \kappa t) \int_0^t d\tau \frac{\sqrt{\kappa\tau}}{t-\tau}I(\tau)
\\
&\approx&
\eta_0(x,t)
\frac{1}{4}
\ln (t)
\eeqa




 Step 11 — RG improvement

Write:

\beq
\log t
=
\log(t/\tau)+\log\tau
\eeq

Absorb:

\beq
g a_1\log\tau
\eeq

into amplitude.

RG condition:

\beq
\tau\frac{dA}{d\tau}
=
g a_1 A
\eeq

Solution:

\beq
A(t)=t^{g/2}
\eeq

Final long-time scaling:

\beq
\eta(x,t)
\sim
t^{-1/2+g/2}
f!\left(\frac{x}{\sqrt{t}}\right)
\eeq

---

 The real conceptual punchline

The logarithm came from:

\beq
\int_0^t \frac{d\tau}{\tau}
\eeq

and that, in turn, came from:

\beq
\frac{\chi}{\tau}
\exp\left(\frac{-\chi^2}{2\kappa \tau}\right)
\eeq

after rescaling.

So the deep mechanism is:

**diffusive spreading + interface discontinuity → 
scale-free $1/\tau$ memory → logarithm**

That is the RG seed.

---

 Why this mirrors QFT exactly

This integral:

\beq
\int_0^t \frac{d\tau}{\tau}
\eeq

is the time analogue of:

\beq
\int \frac{dq}{q}
\eeq

in momentum-space RG.

That parallel is not accidental—it is the real
 reason RG works here.

\section{finding self-similar function}
Good—this is a very nice place where you can
 **see the interface explicitly** in the
  scaling function, and verify that:

* $\eta(x,t)$ is continuous at $x=0$
* $\partial_x \eta$ has a kink at $x=0$
* flux $D,\partial_x \eta$ is continuous

We'll build the scaling solution and 
derive the matching condition.

---

 1. Look for a similarity solution

For long times, diffusion problems tend 
toward **self-similar form**, so try:

\beq
\eta(x,t)
=
t^{-\alpha}
f(\xi),
\quad
\xi=\frac{x}{\sqrt{t}}
\eeq

From dimensional reasoning (or normalization conservation):

\beq
\alpha=\frac12
\eeq

So:

\beq
\eta(x,t)
=
\frac{1}{\sqrt{t}}
f(\xi)
\eeq

---

 2. Plug into the diffusion equation

We have:

\beq
\partial_t \eta
=
\partial_x!\left(D(x)\partial_x \eta\right)
\eeq

Compute derivatives.

First:

\beq
\partial_x \eta
=
\frac{1}{t}
f'(\xi)
\eeq

Then:

\beq
\partial_x^2 \eta
=
\frac{1}{t^{3/2}}
f''(\xi)
\eeq

Time derivative:

\beq
\partial_t \eta
=
-\frac{1}{2t^{3/2}}
\left[
f + \xi f'
\right]
\eeq

So the PDE becomes:

\beq
-\frac12(f+\xi f')
=
D(\xi),f''
\eeq

where:

\beq
D(\xi)=
\begin{cases}
D_1,& \xi>0\\
D_2,& \xi<0
\end{cases}
\eeq

---

 3. ODE on each side

So we solve:

For $\xi>0$:

\beq
D_1 f''+\frac12(f+\xi f')=0
\eeq

For $\xi<0$:

\beq
D_2 f''+\frac12(f+\xi f')=0
\eeq

---

 4. Solve the ODE

These equations are solved by Gaussians.

General solution:

For $\xi>0$:

\beq
f_+(\xi)
=
A_+
\exp!\left(-\frac{\xi^2}{4D_1}\right)
\eeq

For $\xi<0$:

\beq
f_-(\xi)
=
A_-
\exp!\left(-\frac{\xi^2}{4D_2}\right)
\eeq

So the two halves have **different widths**.

---

 5. Continuity condition at $\xi=0$

We must have:

\beq
f_+(0)=f_-(0)
\eeq

So:

\beq
A_+=A_- \equiv A
\eeq

Thus:

\beq
f(0)=A
\eeq

So the function is continuous.

---

 6. Flux continuity condition

Recall:

\beq
D_1\partial_x \eta(0^+)
=
D_2\partial_x \eta(0^-)
\eeq

Compute derivatives.

Since:

\beq
\partial_x \eta
=
\frac{1}{t}
f'(\xi)
\eeq

we get:

\beq
D_1 f'_+(0)
=
D_2 f'_-(0)
\eeq

---

 7. Compute derivatives

Differentiate:

\beq
f'_+(\xi)
=
-\frac{\xi}{2D_1}
A
e^{-\xi^2/4D_1}
\eeq

So:

\beq
f'_+(0)=0
\eeq

Similarly:

\beq
f'_-(0)=0
\eeq

So leading-order Gaussian gives **zero slope on both sides**.

That means:

Flux continuity is automatically satisfied at leading order.

But:

 the **next-order correction** generates the kink.

That's where the discontinuity truly appears.

---

 8. Include next-order correction

Because diffusivities differ, the full solution is:

\beq
f(\xi)
=
A
e^{-\xi^2/4D_i}
\left[
1+\epsilon g(\xi)
\right]
\eeq

Now:

\beq
f'(0^+)\neq f'(0^-)
\eeq

while:

\beq
D_1 f'(0^+)=D_2 f'(0^-)
\eeq

This produces the **kink**.

---

 9. What the kink looks like

Graphically:



Key features:

* Continuous at $\xi=0$
* Different curvature left/right
* Slope changes abruptly

That slope change encodes:

\beq
D_1\neq D_2
\eeq

---

 10. Physical interpretation

If:

\beq
D_1>D_2
\eeq

then:

* right side diffuses faster
* distribution spreads wider to the right
* gradient adjusts to keep flux continuous

So:

* profile is skewed
* center remains continuous

---

 11. Where RG modifies this

The RG calculation earlier changes:

\beq
\eta(x,t)
=
t^{-1/2}
\to
t^{-1/2+\epsilon/2}
\eeq

So the final similarity solution becomes:

\beq
\eta(x,t)
=
t^{-1/2+\epsilon/2}
f(\xi)
\eeq

That exponent shift is the 
**anomalous scaling** produced by the interface.

---

 The big conceptual takeaway

The discontinuity in diffusion coefficient shows up as:

* **continuity of the solution**
* **kink in derivative**
* **logarithmic corrections**
* **RG-modified decay exponent**

All four are different faces of the same underlying structure.

---

 A very nice optional next step

If you're interested, the most revealing thing
 to compute next would be:

**the exact ratio**

\beq
\frac{f'(0^+)}{f'(0^-)}
=
\frac{D_2}{D_1}
\eeq

including first-order corrections.

That calculation directly links:

* interface physics
* RG logarithms
* matching conditions

into one unified picture.


\section{Solving Diffusion Equation via Green Functions}


\beq
\call = \partial_t - D \partial^2_x
\eeq

\beq
\call \eta(x,t) = S(x,t)
\eeq

\begin{claim}
The solution to the equation
\beq
\call \eta (t,x) = S(t,x)\quad
\eeq
for $x>0$ and $t>0$
with initial conditions
\beq
\eta(0, x)= \eta_\rvx(x)
\eeq
and Dirichlet boundary conditions:

\beq
\eta(t, 0) = \eta_\rvt(t)
\eeq
is

\beqa
\eta(t, x)&=&
\int_0^{t+} d\tau\;
\int_0^\infty d\chi\;
G(x,t;\chi,\tau)
f(\chi,\tau)
\\
&&+
\int_0^{\infty} d\chi\;
G(x,t;\chi, 0)
 \eta_{\rvx}(\chi)
\\
&&-D
\int_0^{t+} d\tau\;
[\partial_\chi G(x,t;\chi,\tau)]_{\chi=0}
 \eta_\rvt(\tau)
\eeqa
where $G(x,t;\chi,\tau)$ satisfies 
the following equations for $x>0$ and $t>0$


\beq
\call G(x,t;\chi, \tau)=\delta(x-\chi)\delta(t-\tau)
\eeq
and

\beq
G(x=0,t;\chi,\tau)=0
\eeq
\end{claim}
\proof


\beq
\call= \partial_t -D \partial_x^2,\quad
\call G(x,t;\chi, \tau)=\delta(x-\chi)\delta(t-\tau)
\eeq


\beq
\call^\dagger= -\partial_t -D \partial_x^2,
\quad \call^\dagger G'(x,t;\chi, \tau)=\delta(x-\chi)\delta(t-\tau)
\eeq



\beq
P_\rvx = \int_0^\infty d\chi\; \ket{\rvx=\chi}\bra{\rvx=\chi}
\eeq


\beq
P_\rvt = \int_0^t d\tau\; \ket{\rvt=\tau}\bra{\rvt=\tau}
\eeq

\newcommand{\kett}[2]{\left|\begin{array}{c}1\\2\end{array}\right\rangle}

\newcommand{\brat}[2]{\left\langle\begin{array}{c}1\\2\end{array}\right|}

\newcommand{\stackt}[2]{
\left\{\begin{array}{c}1\\2\end{array}\right\}}

\newcommand{\cev}[1]{\reflectbox{\ensuremath{\vec{\reflectbox{\ensuremath{1}}}}}}

\beq
\partial_{\rvx} = 
\int_0^\infty d\chi\; \ket{\rvx=\chi}\partial_\chi\bra{\rvx=\chi}
\eeq

\beq
\partial_{\rvt} = 
\int_0^\infty dt\; \ket{\rvt=\tau}\partial_\tau\bra{\rvt=\tau}
\eeq

\beq
\rv{\call}= \partial_\rvt-D\partial^2_\rvx
\eeq

$f(x)\cev{\partial_x} a(x)= a(x) \partial_x f(x)$

$\stackt{A}{B}= A\otimes B$

$\kett{x}{t}=\ket{x}\otimes\ket{t}$

\beq
\boxed{\brat{x}{t}\ket{\eta}=
\brat{x}{t}G
\stackt{P_\rvx}{P_\rvt}\ket{S}
+
\brat{x}{t}G\stackt{P_\rvx}{P_0}\kett{\eta_\rvx}{0}
-D
\brat{x}{t}G\stackt{\cev{\partial}_{\rvx}}{P_\rvt}\kett{0}{\eta_\rvt}}
\eeq

\beq
I = \brat{x}{t}G\rv{\call}
\ket{\eta}-
\bra{\eta}\rv{\call^\dagger}
G'\kett{x}{t}
\eeq

\beq
I_\rvt = 
\brat{x}{t}G
\stackt{1}{\partial_\rvt}
\ket{\eta}-
\bra{\eta}
\stackt{1}{(-\partial_\rvt)}
G'\kett{x}{t}
\eeq

\beq
I_\rvx = 
\brat{x}{t}G\partial_\rvx^2
\ket{\eta}-
\bra{\eta}\partial_\rvx^2
G'\kett{x}{t}
\eeq

\beq
\boxed
{I = I_\rvt -D I_\rvx}
\eeq

\beq
\call^\dagger \brat{x}{t}G\kett{\chi}{\tau
}=\brat{x}{t}\rv{\call}^\dagger G'\kett{\chi}{\tau}=
\stackt{\delta(x-\chi)}{\delta(t-\tau)}
\implies\quad
\rv{\call}^\dagger G' =1
\eeq



\beq
\rv{\call}\ket{\eta}=\ket{S}
\eeq

\beq
I = \brat{x}{t}G\ket{S}-\bra{\eta}\kett{x}{t}
\eeq

\beqa
I_\rvt &=&
\int_0^t d\tau
\brat{x}{t}G
\stackt{1}{\ket{\tau}\partial_\tau\bra{\tau}}
\ket{\eta}
+
\bra{\eta}
\stackt{1}{\ket{\tau}\partial_\tau\bra{\tau}}
G'\kett{x}{t}
\eeqa

\beq
\bra{t}G\ket{\tau}= \bra{\tau}G'\ket{t}
\eeq

\beq
\partial_\tau (\ket{\tau}\bra{\tau})=
\ket{\tau}\partial_\tau\bra{\tau}
+\ket{\tau}\cev{\partial}_\tau\bra{\tau}
\eeq


\beqa
I_\rvt &=&
\int_0^t d\tau
\brat{x}{t}G
\stackt{1}{\partial_\tau (\ket{\tau}\bra{\tau})}
\ket{\eta}
\eeqa

\beq
\int_0^t
d\tau\;\partial_\tau (\ket{\tau}\bra{\tau})
=
\ket{t}\bra{t}-\ket{0}\bra{0}
\eeq


\beq
\boxed{I_\rvt =
\brat{x}{t}G
\stackt{P_\rvx}{\ket{t}\bra{t}-\ket{0}\bra{0}}
\ket{\eta}}
\eeq

\beqa
I_\rvx &=& 
\int_0^\infty d\chi
\brat{x}{t}G
\stackt{\ket{\chi}\partial_\chi^2\bra{\chi}}{1}
\ket{\eta}
-
\bra{\eta}
\stackt{\ket{\chi}\partial_\chi^2\bra{\chi}}{1}
G'\kett{x}{t}
\eeqa

\beq
\ket{\chi}\partial_\chi^2\bra{\chi}
-\ket{\chi}\cev{\partial}_\chi^2\bra{\chi}=
\partial_\chi(\ket{\chi}\partial_\chi\bra{\chi})
-(\ket{\chi}\cev{\partial}_\chi\bra{\chi})\cev{\partial}_\chi
\eeq

\beq
I_\rvx=
\int_0^\infty d\chi
\brat{x}{t}G
\stackt{\partial_\chi(\ket{\chi}\partial_\chi\bra{\chi})
-(\ket{\chi}\cev{\partial}_\chi\bra{\chi})\cev{\partial}_\chi}{1}
\ket{\eta}
\eeq

\beqa
\int_0^\infty d\chi
\partial_\chi(\ket{\chi}\partial_\chi\bra{\chi})
-(\ket{\chi}\cev{\partial}_\chi\bra{\chi})\cev{\partial}_\chi
&=&
[\ket{\chi}\partial_\chi\bra{\chi}
-
\ket{\chi}\cev{\partial}_\chi\bra{\chi}]_0^\infty
\\
&=&
-[\ket{\chi}\partial_\chi\bra{\chi}
-
\ket{\chi}\cev{\partial}_\chi\bra{\chi}]_{\chi=0}
\eeqa

\beq
I_\rvx=
\brat{x}{t}G
\stackt{[-\ket{\chi}\partial_\chi\bra{\chi}
+
\ket{\chi}\cev{\partial}_\chi\bra{\chi}]_{\chi=0}}{1}
\ket{\eta}
\eeq


\beq
\boxed{I_\rvx=\brat{x}{t}G
\stackt{[
\ket{\chi}\cev{\partial}_\chi\bra{\chi}]_{\chi=0}}{P_\rvt}
\ket{\eta}}
\eeq
\qed


\beq
G^\pm(x,t;\chi,\tau)
=
\frac{1}{\sqrt{4\pi D(t-\tau)}}
\exp\left[-\;\frac{(x\pm\chi)^2}{4D(t-\tau)}\right]
\eeq

\beq
G = G^+ - G^-
\eeq
At the boundary
\beq
\left.
\partial_\chi G
(x,t;\chi,\tau)
\right|_{\chi=0}
=
\frac{x}{\sqrt{4\pi D (t-\tau)^3}}
\exp\left(
-\frac{x^2}{4D(t-\tau)}
\right)
\eeq



\end{document}